\subsection{将查询排序以计算表值}

在元数为 $t$ 的站点，与之相接的各比特给出长度为 $2^t$ 的表中某一表项的地址。同时计算该站点的许多副本时，所选表项的乘积只取决于各个地址出现的次数。将查询排序即可求出这些次数，并将查表成本分摊到整个批次。

\begin{lemma}[批量查表]\label{lem:batch-table} 假设定理~\ref{thm:sort} 关于通常宽度的稳定排序界在 $k$ 上成立，指数为 $p>1$。设一个由二元线路与取值于 $K$ 的表组成的网络表示张量 $T$，其中 $K/k$ 为任意域扩张。若站点 $v$ 的元数为 $t_v\le u$，并且 $s=\sum_v(1+t_v)$，则 $T^{\otimes 2^u}$ 有一个精确的三价 $K$-网络，其规模为 \[
                  O\bigl(2^u s(1+u)^{p-1}\bigr).
\] 隐含常数与 $K$ 及各表的表项均无关。\end{lemma}

\begin{proof} 先考虑单张表 $f\colon\{0,1\}^t\to K$，其中 $t\ge1$，并将 $2^t$ 个查询地址 $x_1,\ldots,x_{2^t}$ 作为张量的外露腿。若地址 $z$ 出现 $c_z$ 次，所需乘积即为 $\prod_z f(z)^{c_z}$，其中未出现的地址对应空积。为每个 $z$ 创建一行字典记录，用 $t+1$ 个比特猜测其出现次数。在计数的第 $j$ 位（$0\le j\le t$）放置一个一元因子，其两个分量为 \[
                    (1,\ f(z)^{2^j}).
\] 该因子的位置已经确定了 $z$，因此两个分量都是网络的固定系数。

为每个查询建立单独的一行，并将其中的计数字段置零。将 $2^{t+1}$ 行先按地址、再按行的类型排序，使每一字典行位于对该地址的所有查询之后。计数器从零开始，每遇到一个查询便加一。遇到字典行时，要求计数器与该行携带的计数相符，再将计数器归零。即使一个地址没有查询，其计数也得到核验。对固定的外露地址，计数、排序后的行及计数器状态恰有一种合法取法。因此，在任意特征下，收缩的结果都等于 $\prod_i f(x_i)$，即使某个表项为零也成立。

用逐位进位加法实现计数器需 $O(2^t(1+t))$ 个门。将宽 $O(t+1)$ 比特的各行排序，成本为 $O(2^t(1+t)^p)$。排序网络在 $k$ 上实现一个零一张量，经标量扩张后便在 $K$ 上实现同一张量。

将每个元数为 $t$ 的站点的 $2^u$ 个副本分成 $2^{u-t}$ 批，每批 $2^t$ 个。元数为零的站点直接提供一个标量。在每个副本内依照原网络连接外露线路。所得网络收缩为 $T^{\otimes 2^u}$，其规模至多为下式的常数倍：\[
  2^u\sum_v(1+t_v)^p
       \ \le\ 2^u(1+u)^{p-1}\sum_v(1+t_v).
\] \end{proof}
